\section{Actual homogeneous generators for the raw chart estimate}

\begin{lemma}[Finite homogeneous generators of the same whole ideal]
\label{lem:c0n-finite-homogeneous-whole-generators}
\lean{EnvelopeProperty.ProjSubvariety.MultiplierBaseIdeals.exists_finset_homogeneous_generators}
\leanok
Over a Noetherian coefficient ring, every actual homogeneous ideal in a
polynomial ring with finitely many variables has a finite set of genuine
homogeneous members generating that entire same ideal.
\end{lemma}
\begin{proof}
Every element of a homogeneous ideal is the finite sum of its homogeneous
components, which remain members of the same ideal.  Thus all genuine
homogeneous members span the whole ideal.  Noetherianity selects a finite
subset of this actual spanning set, retaining both membership and
homogeneous degree for every selected generator.
\end{proof}

\begin{lemma}[Whole canonical cone ideals have actual finite homogeneous polynomial generators]
\label{lem:c0n-raw-canonical-cone-generators}
\uses{lem:c0n-finite-homogeneous-whole-generators,lem:c0n-canonical-scheme-cone-representative,lem:c0n-relative-homogeneous-parameter-coefficients}
\lean{EnvelopeProperty.ProjSubvariety.MultiplierBaseIdeals.exists_finset_coneIdeal_homogeneous_generators}
\leanok
For every actual ideal sheaf $K$ of the original parameter family, its
same canonical cone ideal has a finite set $G$ of whole original ambient
polynomials such that every polynomial is homogeneous of one $x$-degree,
its actual parameter exchange belongs to the canonical homogeneous
representative of $K$, and the image of $(G)$ under the original cone
quotient equals the entire canonical cone ideal exactly.
No bound on these degrees is asserted.
\end{lemma}
\begin{proof}
Apply the preceding lemma to the canonical homogeneous representative of
the same whole $K$, over the Noetherian coefficient ring $k[t]$.
Apply the inverse of the actual original parameter exchange to each of
these finitely many homogeneous members.  Coefficient preservation shows
that every resulting whole polynomial has the same homogeneous
$x$-degree in all of its parameter coefficients.
The forward parameter exchange returns the original representative
member, so membership is retained exactly.
Extension of ideals under the inverse parameter exchange followed by
the original cone quotient equals extension under their actual composite,
which is the defining original parameter cone quotient.
The selected members span the entire canonical representative; therefore
their resulting cone ideal is exactly the original canonical cone ideal.
\end{proof}
